\chapter{Conclusion}
\label{ch:concl}

\section{Expected Outcomes}

At the end of the project, the following outcomes are expected. They are goals
for the implementation phase, not results that have already been achieved.

\begin{itemize}
  \item \textbf{A working framework.} A Python-based PurpleAgents framework with
  red and blue agents, a dual-reasoning engine, reflection, validation, an
  auditable memory, and human approval, operated through a CLI.
  \item \textbf{A measurable improvement in defense.} Over repeated
  attack--defend rounds, the detection rate is expected to rise and the attack
  success rate to fall, without a human writing the fixes.
  \item \textbf{Fewer wrong conclusions.} Dual-path reasoning is expected to
  catch judging mistakes that a single path would miss, with uncertain cases sent
  to a human instead of being decided silently.
  \item \textbf{A complete audit trail.} Every attack, decision, fix, and approval
  is expected to be traceable in the memory.
  \item \textbf{Clear CLI reports.} Command-line summaries that show each
  round's results, the trend in detection rate, and the lessons learned.
\end{itemize}

\section{Conclusion}

Phase~1 of PurpleAgents has established a clear design and a feasibility
foundation for an autonomous purple-team framework. The framework continuously
attacks and defends an AI model and learns from every missed attack.

A verified research gap was identified through a review of recent work. Existing
systems handle attack, defense, or reasoning separately, but none of them
integrate all three into one closed, self-improving loop.

The proposed design fills this gap. It uses a dual-reasoning engine that accepts
a conclusion only when the evidence-first and hypothesis-first paths agree. It
adds a reflection step that turns misses into fixes, an auditable memory that
records every decision, and a closed learning loop that runs under human
approval.

The feasibility study confirms that the project can be built with free,
open-source tools on modest hardware. The planned evaluation --- repeated
attack--defend cycles that measure the rise in the defense catch-rate --- will
test the system's core claim: that the defense improves automatically over time.
The remaining work is implementation, testing, and the measurement of real
results in the next phase.
